Misalkan $Y$ merupakan gabungan saling lepas dari semua $U_i$. Pada $Y$ kita definisikan sebuah
\definitionsverweis {relasi ekuivalensi}{}{}
$\sim$ dengan menyatakan titik-titik

\mavergleichskette
{\vergleichskette
{ x_i
}
{ \in }{ U_i
}
{  }{
}
{    }{
}
{   }{
}
}
{}{}{}
dan

\mavergleichskette
{\vergleichskette
{ x_j
}
{ \in }{ U_j
}
{  }{
}
{    }{
}
{   }{
}
}
{}{}{}
ekuivalen jika

\mavergleichskette
{\vergleichskette
{ x_i
}
{ \in }{ U_{ij}
}
{  }{
}
{    }{
}
{   }{
}
}
{}{}{,}

\mavergleichskette
{\vergleichskette
{ x_j
}
{ \in }{ U_{ji}
}
{  }{
}
{    }{
}
{   }{
}
}
{}{}{}
dan

\mavergleichskette
{\vergleichskette
{ \varphi_{ji} (x_i)
}
{ = }{ x_j
}
{  }{
}
{    }{
}
{   }{
}
}
{}{}{}.
Sifat-sifat relasi ekuivalensi dijamin oleh syarat koklus; lihat
soal tentang relasi ekuivalensi.
Kita tetapkan

\mavergleichskettedisp
{\vergleichskette
{X
}
{ \defeq} { Y/ \sim
}
{ } {
}
{ } {
}
{ } {
}
}
{}{}{}
dan melengkapi $X$ dengan
\definitionsverweis {topologi hasil bagi}{}{.}
Komposisi-komposisi

\mathdisp {U_i \longrightarrow Y \longrightarrow X} {  }

adalah $\psi_i$, dan $V_i$ merupakan citra pemetaan-pemetaan ini. Jadi diperoleh homeomorfisme-homeomorfisme
\maabb {\psi_i} {U_i } { V_i
} {}.
Untuk

\mavergleichskette
{\vergleichskette
{ x
}
{ \in }{ U_i
}
{  }{
}
{    }{
}
{   }{
}
}
{}{}{}

\mavergleichskettedisp
{\vergleichskette
{ \psi_i (x)
}
{ \in} { V_j
}
{ } {
}
{ } {
}
{ } {
}
}
{}{}{}
berlaku jika dan hanya jika

\mavergleichskette
{\vergleichskette
{ x
}
{ \in }{ U_{ij}
}
{  }{
}
{    }{
}
{   }{
}
}
{}{}{},
sebab tepat dalam kasus ini $x$ diidentifikasi dengan
\mathl{\varphi_{ji}(x)}{}.
Oleh karena itu,

\mavergleichskette
{\vergleichskette
{ \psi_i(U_{ij})
}
{ = }{ V_i \cap V_j
}
{  }{
}
{    }{
}
{   }{
}
}
{}{}{.}
Komutativitas diagram

\mathdisp {\begin{matrix} U_{ij}   & \stackrel{ \varphi_{ji} }{\longrightarrow} &  U_{ji}   & \\ & \!\!\! \!\! \psi_i \searrow & \downarrow \psi_j \!\!\! \!\! & \\ & &  V_i \cap V_j   & \!\!\!\! \!\!\!\!  \\ \end{matrix}} {  }

mengikuti dengan cara yang sama.
